\chapter{Open questions}
\label{ch:open}

\newthought{No-one has won} the long-context architecture argument. This
chapter is a short, opinionated list of the questions whose answers will
decide it over the next few years.

\section{Selection algorithms at frontier scale}

For content-dependent sparse attention to live in the bottom-right
quadrant of Chapter~\ref{ch:quadrant}, the selection function has to be
fast \emph{and} accurate \emph{and} differentiable. We have plenty of
algorithms that hit two of those: ANN indexes are fast and accurate but
not differentiable; learned routers are differentiable and fast but lose
accuracy; hierarchical clustering is differentiable and accurate but slow
to update during training.

The open question is whether SSA-style mechanisms can keep their
selection function quality at the 1B-parameter / 1T-token scale where the
real customers live. SubQ's published results suggest yes; independent
reproduction is the next move.

\section{Training infrastructure for the bottom-right quadrant}

A standard transformer trains in a well-understood, GPU-friendly way:
matrix multiplies dominate, FlashAttention handles the rest, communication
patterns are predictable. A content-dependent sparse model breaks every
one of those assumptions. The selection function makes computation
data-dependent, the sparsity makes communication patterns irregular, and
the routing produces load imbalance across devices.

How to train these models at scale, on cheap commodity hardware, with
mixed precision and tensor parallelism, is genuinely an open engineering
problem. The architectural papers tend to be quiet on this; the
implementations that work are mostly proprietary.

\section{Evaluation that survives Goodhart}

Every benchmark in the long-context space has been gamed within months
of release.\sidenote{Original needle-in-a-haystack tests were trivial
once everyone trained on synthetic NIAH data; the move to RULER and now
MRCR v2 is the response. Each new benchmark eventually stops
discriminating, and a harder one has to be built.} The community's
ability to come up with new benchmarks faster than vendors can train to
them is the only thing keeping this honest. Expect MRCR v3, v4, etc.

\section{Inference economics, not training economics}

Frontier model training has consumed the field's budget for five years.
The question that decides 2027--2030 is whether anyone can serve a
high-quality 1M-token model for under a cent per million input tokens.
The arithmetic doesn't work for dense transformers; whether it works for
SSA, Mamba-hybrids, or some still-unnamed architecture is genuinely
unsettled. Whoever wins the inference economics question wins the
agentic-applications business.

\keyidea{The next five years of LLM research are about inference, not
training. Whoever can serve frontier-quality 1M-token context cheaply
wins the agent platform business.}

\begin{itemize}
\item \textbf{Data-independent hashing} (LSH, random projections): cheap and
differentiable-free, but buckets may miss true nearest neighbours unless bit
width is large---hard when keys change every layer.
\item \textbf{Data-dependent clustering} (k-means / IVF codebooks): amortises
routing lookups, but centroids must track moving keys and degrade under
distribution shift.
\item \textbf{Learned routers} (small MLPs, gating scores): flexible and
trainable end-to-end, but risk collapsing diversity if routing saturates (see
also DeepSeek-style sparsity \citep{deepseek2024dsa}).
\item \textbf{Hierarchical cluster-then-attend}: coarse stage prunes most keys,
fine stage reranks candidates; good recall but two-pass latency.
\item \textbf{Hybrid retrieve-and-rerank}: ANN index proposes candidates, dense
dot-products rerank top-\(m\); closest to production retrieval stacks, still
needs careful batching inside autoregressive decoding.
\end{itemize}

\begin{table}[htbp]
\centering
\footnotesize
\begin{tabular}{@{}lcc@{}}
\toprule
Architecture (same 70B-active class) & \$/M input @ 1M ctx (illustrative) & Notes \\
\midrule
Dense + FA-2 & \(\$\,30\)--\(\$\,60\) & Quadratic prefill + full KV dominate. \\
GQA + sliding window & \(\$\,8\)--\(\$\,20\) & Linear-ish layers shrink active pairs; KV still \(\mathcal{O}(n)\). \\
Hybrid Mamba--Attn (Jamba-class) & \(\$\,5\)--\(\$\,12\) & Few attention layers cut peak FLOPs; some full KV blocks remain. \\
SSA / SubQ-style routing & \(\$\,1\)--\(\$\,4\) & Uses vendor-reported \(\sim 50\times\) prefill win vs.\ FA-2 at 1M as anchor \citep{subq2026ssa}; real cost depends on SLO. \\
\bottomrule
\end{tabular}
\caption{Order-of-magnitude USD per million input tokens at \(n=10^{6}\),
assuming mixed B200/H100 fleets, utilisation \(\approx 0.2\)--\(0.4\), and
blended rent of a few USD per GPU-hour---meaningful only for comparing ratios
between rows.}
\label{tab:dollar-per-m}
\end{table}

\noindent
Absolute dollars are not knowable in print; the point of the table is relative
ranking under identical accounting rules.
